\section{Introduction}\label{sec:intro}

Let $\A$ be a finite alphabet and $\sh$ the left shift on $\A^{\Z}$. By the Curtis--Hedlund--Lyndon theorem \cite{Hedlund1969} the endomorphisms of the full shift are the block maps, and its automorphisms the invertible block maps. The group $\Aut(\A^{\Z})$ is enormous: it contains every finite group and free groups \cite{BLR1988}, and every countable, locally finite, residually finite group \cite{KimRoush1990}, and its center consists of the powers of $\sh$ \cite{Ryan1972}. A reversible cellular automaton $F\in\Aut(\A^{\Z})$ determines a two-dimensional subshift, its \emph{spacetime} $X_F\subset\A^{\Z^2}$, whose rows are the configurations $F^n(c)$, $n\in\Z$ (see Cyr, Franks and Kra \cite{CFK2019} for the spacetime of a general endomorphism). The automorphism group of $X_F$ is the centralizer of $F$ in $\Aut(\A^{\Z})$ (Lemma~\ref{lem:spacetime}), and it always contains $\langle\sh,F\rangle$. This paper asks, for a natural finite family of nonlinear reversible automata, whether it contains anything else.

The family consists of the second-order elementary rules, built by the second-order technique of Fredkin \cite{Margolus1984,TM1990,Wolfram2002}; as models of statistical mechanics they were introduced by Takesue \cite{Takesue1987,Takesue1989}, who reduced them to the $88$ classes considered here \cite[Table~I]{Takesue1989}, determined their additive conserved quantities with densities on two adjacent cells \cite[Table~II]{Takesue1989}, determined, for words of length at most $6$, which classes have a local conservation law, and exhibited walls in $40$ of them \cite[Section~4, Fig.~2, Table~III]{Takesue1989} (see also \cite{HT1991,Takesue1995}). For an elementary rule $f\colon\{0,1\}^3\to\{0,1\}$ define, on $\A=\{0,1\}^2$ with cells written $(x,y)$,
\begin{equation}\label{eq:F}
F_f(c)_i=\bigl(y_i,\;x_i\oplus f(y_{i-1},y_i,y_{i+1})\bigr).
\end{equation}
Reading $y$ as the current row and $x$ as the previous one, this is the recursion $y_{t+1}=y_{t-1}\oplus f(y_t)$, reversible for every $f$, with $F_f^{-1}(c)_i=\bigl(y_i\oplus f(x_{i-1},x_i,x_{i+1}),\,x_i\bigr)$. Following Takesue \cite{Takesue1987}, one writes $f$R for $F_f$; Wolfram discusses $37$R as a reversible automaton with complex behaviour \cite[pp.~452--457]{Wolfram2002}. Reflection and the cellwise complement conjugate $F_f$ to $F_{f'}$ for the reflected and complemented rules and preserve everything we consider (Lemma~\ref{lem:classes}), so the question is one about $88$ classes, the classes of Takesue \cite[Table~I]{Takesue1989}.

\subsection{The classification}

We assign to each class one of five statuses (Definition~\ref{cl:def:status}): (G), $\Aut(X_{F_f})$ contains $\bigoplus_{m\ge1}S_m$ and hence every finite group; (Z), $\Aut(X_{F_f})=\langle\sh,F_f\rangle\cong\Z^2$; (S), $\Aut(X_{F_f})=\langle\sh,F_f\rangle\times\langle\kappa\rangle\cong\Z^2\times\Z/2$, where $\kappa$ is the cellwise complement; (A), $f$ is affine and $\Aut(X_{F_f})$ is a free abelian group strictly larger than $\langle\sh,F_f\rangle$; and (O), open.

\begin{theorem}[Classification]\label{thm:main}
The $88$ classes of second-order elementary automata, and the $256$ rules, are distributed among the five statuses as follows.
\begin{center}
\begin{tabular}{lrr}
\toprule
status & classes & rules\\
\midrule
(G) every finite group & $43$ & $104$\\
(Z) $\Aut=\langle\sh,F\rangle\cong\Z^2$ & $11$ & $40$\\
(S) $\Aut=\langle\sh,F\rangle\times\langle\kappa\rangle\cong\Z^2\times\Z/2$ & $4$ & $8$\\
(A) affine, $\Aut\cong\Z^2$ or $\Z^3$, $\supsetneq\langle\sh,F\rangle$ & $2$ & $6$\\
(O) open & $28$ & $98$\\
\midrule
total & $88$ & $256$\\
\bottomrule
\end{tabular}
\end{center}
The classes of each status are those of Tables~\ref{cl:tab:finite}, \ref{cl:tab:decided} and~\ref{cl:tab:open}. For each of the $28$ open classes, every automorphism of $X_{F_f}$ of radius at most $3$ is of the form $\sh^aF_f^b$ (computer-assisted), and $(a,b)\mapsto\sh^aF_f^b$ is an isomorphism $\Z^2\to\langle\sh,F_f\rangle$. For the class of $15$, $\Aut(X_{F_f})\cong\Z^2$ contains $\langle\sh,F_f\rangle$ with index $2$; for the class of $105$, $\Aut(X_{F_f})\cong\Z^3$; and for the class of $95$, $\Aut(X_{F_f})\cong(\mathcal N\times\mathcal N\times\Z)/\langle(\sh^{-1},\sh,-2)\rangle$, where $\mathcal N$ is the group of automorphisms of the golden mean shift that preserve the number of $1$s of every finite point \cite{PaperE}.
\end{theorem}

The theorem is proved in Section~\ref{cl:sec:main}, where the tables list the classes with the result that settles each entry. The statuses (G), (Z), (S) and (A) are quoted from the companion papers \cite{PaperA,PaperB,PaperC,PaperE} (Section~\ref{sec:imported}); what is proved here is the assembly and the statement on the open classes. The status (O) records only that none of the other four is known; its mathematical content is the census of Theorem~\ref{cl:thm:census} together with Theorem~\ref{cl:thm:A}. Whether $\Aut(X_{F_f})=\langle\sh,F_f\rangle$ for every open class is Question~\ref{q:open}. The censuses do not decide it: for the class of $95$ they contain no automorphism outside $\langle\sh,F_f\rangle$, and the group contains every finite group \eqref{imp:95}. They do show that for every open class the block maps of radius at most $4$ that commute with $F_f$ are the maps $\sh^aF^b$ and the constants, whereas for $95$ there are $20$ further ones, none of them injective, at radius $3$.

\subsection{Results used from the companion papers}\label{sec:imported}

The statuses (G), (Z), (S) and (A), and the isomorphism $\langle\sh,F_f\rangle\cong\Z^2$, rest on the following results, which are proved in four companion papers \cite{PaperA,PaperB,PaperC,PaperE}; none of them is proved here. We state them in the form in which they are used.

\begin{enumerate}
\renewcommand{\theenumi}{I\arabic{enumi}}
\renewcommand{\labelenumi}{(\theenumi)}
\item\label{imp:walls} \emph{Fixed walls} \cite{PaperC}. Let $F$ be a reversible cellular automaton of radius $1$ on $\A^{\Z}$ and $W$ a nonempty word such that $F(c)$ contains $W$ at every position at which $c$ does. If at least two symbols $s$ have the property that $W\,s\,W$ contains $W$ only at its two ends, then $\bigoplus_{m\ge1}S_m$ embeds in $\Aut(X_F)$. Of the $256$ elementary rules, $81$, in $31$ classes, have such a word of length at most $3$ (a finite check); these words are the fixed walls (i)--(viii) of Takesue \cite[Fig.~2, Table~III]{Takesue1989}.
\item\label{imp:54} \emph{Rule $54$} \cite{PaperC}. For every $m\ge1$ there is an injective homomorphism $\Psi_m\colon S_m\to\Aut(X_{F_{54}})$ whose image consists of block maps of radius at most $3m+3$, and the images for different $m$ commute; so $\bigoplus_{m\ge1}S_m$ embeds in $\Aut(X_{F_{54}})$.
\item\label{imp:133} \emph{A wall of length $4$} \cite{PaperC}. The word $0330$ is fixed by $F_{133}$ in every context (it is the wall (ix) of \cite[Table~III and p.~391]{Takesue1989}), and \eqref{imp:walls} applies to the class $\{133,161\}$, with automorphisms of radius $4$ outside $\langle\sh,F\rangle$. Every word fixed in every context by some $F_f$ lies in $\{0,3\}^k$, and among the $57$ classes with no such word of length at most $3$, the class of $133$ is the only one with such a word of any length (the second statement by a reduction to a finite search, computer-assisted in \cite{PaperC}).
\item\label{imp:hidden} \emph{Hidden-bit and time-periodic walls} \cite{PaperC}. For the classes of $59$ and $27$ (walls carrying an internal bit invisible to their neighbours) and of $127$, $179$ and $223$ (walls fixed by a power of $F$ in every context), $\bigoplus_{m\ge1}S_m$ embeds in $\Aut(X_{F_f})$. A computer search among the closures of single words under $F^{\pm1}$ (of at most $300$ words), with windows of length at most $6$ and a hidden involution satisfying the conditions (W1) and (H1) of \cite{PaperC}, in the $57$ classes without a fixed word of length at most $3$, finds such families only in these five classes \cite{PaperC}. The walls of these five classes are built from the periodic walls of \cite[Section~4]{Takesue1989}.
\item\label{imp:rot} \emph{Rotating walls} \cite{PaperC}. For $f\in\{126,219\}$, runs of cells that $F_f$ rotates through the alphabet in every context give an embedding of $\bigoplus_{m\ge1}S_m$ in $\Aut(X_{F_f})$; these runs consist of the walls (xv) and (xvi) of \cite[Fig.~2, Table~III]{Takesue1989}.
\item\label{imp:90} \emph{Markers for $90$ and $165$} \cite{PaperC}. For $f\in\{90,165\}$, $F_f$ has no blocking word (Section~\ref{sec:blocking}), and $(\bigoplus_{m\ge1}S_m)^2$ embeds in $\Aut(X_{F_f})$. The automorphisms of $X_{F_{90}}$ of radius at most $2$ are affine (computer-assisted in \cite{PaperC}; the enumeration is repeated in Table~\ref{tab:census}).
\item\label{imp:affine} \emph{Affine classes} \cite{PaperC}. For the classes of $15$, $60$, $105$ and $153$, every automorphism of $X_{F_f}$ is affine, by the isomorphism rigidity theorem of Kitchens and Schmidt \cite{KS2000}, and: for $f\in\{170,240\}$, $\Aut(X_{F_f})=\langle F,\mathrm{id}+F\rangle\cong\Z^2$, with $(\mathrm{id}+F)^2=\sh^{\pm1}F$, so that $\langle\sh,F\rangle$ has index $2$, and for $f\in\{15,85\}$ the same holds with the generator $c\mapsto c+F(c)+(1,1)^{\Z}$; for the classes of $60$ and $153$, $\Aut(X_{F_f})=\langle\sh,F\rangle\times\langle\kappa\rangle\cong\Z^2\times\Z/2$; for $f=150$, $\Aut(X_{F_f})=\langle\sh,F,\mathrm{id}+\sh F\rangle\cong\Z^3$, and for $f=105$ the same holds with $c\mapsto c+\sh F(c)+(1,1)^{\Z}$ in place of $\mathrm{id}+\sh F$. For $f=255$, $F$ adds $1$ modulo $4$ in every cell after a relabelling of the cells by $\Z/4$, $\Aut(X_{F_{255}})\cong\Z/4\times\mathcal G_4$, where $\mathcal G_4$ is the group of automorphisms $\psi$ of the full shift on $\Z/4$ with $\psi-\mathrm{id}=(\sh-\mathrm{id})J$ for a block map $J$, and $\bigoplus_{m\ge1}S_m$ embeds in $\mathcal G_4$.
\item\label{imp:perm} \emph{Permutive rules} \cite{PaperB}. If $f$ is not affine and is permutive in its left or its right variable, then $\Aut(X_{F_f})=\langle\sh,F_f\rangle\cong\Z^2$. These are the $16$ rules of the classes of $30$, $45$, $75$ and $135$.
\item\label{imp:half} \emph{Rules that read two neighbours} \cite{PaperB}. For the $8$ rules of the classes of $63$ and $187$, $\Aut(X_{F_f})=\langle\sh,F_f\rangle\cong\Z^2$.
\item\label{imp:110} \emph{Rules $110$, $14$, $46$, $79$, $24$ and $189$} \cite{PaperA}. For the $14$ rules of the classes of $110$, $14$, $46$ and $79$, $\Aut(X_{F_f})=\langle\sh,F_f\rangle\cong\Z^2$. For the $4$ rules of the classes of $24$ and $189$, $\Aut(X_{F_f})=\langle\sh,F_f\rangle\times\langle\kappa\rangle\cong\Z^2\times\Z/2$, where $\kappa$ is the cellwise complement \cite[Theorem~1.1]{PaperA}. All proofs in \cite{PaperA}, including those for $24$ and $189$, are by hand.
\item\label{imp:order} \emph{Infinite order} \cite{PaperC}. If $f$ is not a function of its middle argument alone, then $(a,b)\mapsto\sh^aF_f^b$ is injective, so $\langle\sh,F_f\rangle\cong\Z^2$ \cite[Theorem~2.4]{PaperC}. The statement is recalled as Theorem~\ref{cl:thm:A}.
\item\label{imp:94} \emph{Rule $94$} \cite{PaperE}. For $f\in\{94,122\}$, $\Aut(X_{F_f})=\langle\sh,F_f\rangle\cong\Z^2$ \cite[Theorem~A; Theorem~5.9]{PaperE}.
\item\label{imp:95} \emph{Rule $95$} \cite{PaperE}. For $f\in\{95,250\}$, $\Aut(X_{F_f})\cong(\mathcal N\times\mathcal N\times\Z)/\langle(\sh^{-1},\sh,-2)\rangle$, where $\mathcal N$ is the group of automorphisms of the golden mean shift that preserve the number of $1$s of every finite point, and $\bigoplus_{m\ge1}S_m$ embeds in $\mathcal N$, so that $\Aut(X_{F_f})$ contains every finite group; the automorphism group of the spacetime of the reversible automaton known as Rule~54, which is the system on one sublattice of a diagram of $F_{95}$, is $\mathcal N\times\mathcal N$; $F_f$ has no blocking word, and the window graphs $W^{(k,p)}_M$ of Section~\ref{sec:blocking} are strongly connected for all $(k,p)$ and $M$ \cite[Theorem~B; Theorem~6.9, Corollary~6.10, Propositions~6.17 and~6.18]{PaperE}. The smallest radius of an automorphism outside $\langle\sh,F_f\rangle$ is $7$; and the block maps of radius at most $4$ that commute with $F_{95}$, other than the maps $\sh^aF^b$ and one constant map, number $20$ at radius $3$ and $220$ at radius $4$, are induced by pairs of block maps of the golden mean shift that conserve the number of $1$s of finite points and are not both injective, and none of them is injective (computer-assisted) \cite[Proposition~6.15, Lemma~6.20 and Proposition~6.21]{PaperE}.
\item\label{imp:L} \emph{Lemma~L; rules $91$ and $155$} \cite{PaperE}. For $f\in\{91,94,95\}$, a block map of any radius that commutes with $F_f$ and fixes every fixed point of $\sh F_f$ and of $\sh^{-1}F_f$ is the identity \cite[Theorem~C; Theorem~3.1, Corollary~3.9]{PaperE}. For $91$ the group is not determined there; for $155$ the method of proof fails \cite[Proposition~8.1]{PaperE}, and the statement is verified for radius at most $6$ (computer-assisted) \cite[Remark~8.2]{PaperE}.
\end{enumerate}

The proofs of \eqref{imp:perm}--\eqref{imp:110} and of \eqref{imp:94} share a pattern: an automorphism is normalised by an element of $\langle\sh,F\rangle$ to one that fixes certain distinguished structures, and such a map is shown to be the identity. The structures are the light-speed fronts of differences for the permutive rules, and the rightmost difference between two configurations, which is stationary under $F$, for the classes of $63$ and $187$ \cite{PaperB}; the kinks among the fixed points of $\sh^{\pm1}F$ for $110$ and $46$, together with those of $F$ for $14$ and of $\sh^{\pm1}F^3$ for $79$, and for $24$ and $189$ a factor map onto rule $46$ \cite{PaperA}; and for the class of $94$ the kinks among the fixed points of $\sh^{\pm1}F$, which form a shift of finite type of positive entropy, where the scattering of kinks carries the normalisation to all these fixed points and \eqref{imp:L} concludes \cite{PaperE}. For $95$ the first step and \eqref{imp:L} hold as well, but a normalised automorphism need not be the identity; for the open class of $91$, \eqref{imp:L} gives part of the second step.

\subsection{The results of this paper}

What is proved in this paper is the following; everything else in the classification is quoted from \eqref{imp:walls}--\eqref{imp:L}.

\emph{The classification} (Theorem~\ref{thm:main}, Tables~\ref{cl:tab:finite}--\ref{cl:tab:open}). The assembly of the statuses from the imported results: which of them applies to which class, the counts, and the complete list of the permutations of the four cells that commute with $F_f$ (Remark~\ref{cl:rem:kappa}).

\emph{Censuses} (Theorems~\ref{thm:110} and~\ref{cl:thm:census}; computer-assisted). For a fixed radius $R$, commutation with $F_f$ is a finite set of constraints on the $2\cdot4^{2R+1}$ bits of the table of a block map. We enumerate all solutions with a SAT solver, re-verify each by direct evaluation, decide invertibility exactly, and repeat the enumeration with other solvers (Section~\ref{sec:method}). For the open classes and the class of $95$ the commuting block maps of radius at most $3$ are the maps $\sh^aF^b$, the constants, and, for $95$ only, $20$ non-injective maps (Theorem~\ref{cl:thm:census}). For the rules $110$ and $30$ the block maps of radius at most $4$ that commute with $F_f$ are the maps $\sh^aF^b$ and two constant maps (Theorem~\ref{thm:110}). A census at radius $4$ for the open classes, searches for walls, eigenfields and conserved quantities, and the window graphs of all classes are reported in Section~\ref{cl:sec:holdouts} (computer-assisted).

\emph{No walls in $110$R} (Proposition~\ref{cl:prop:nowalls}). Every wall mechanism of \cite{PaperC} needs a cell next to the permuted information whose trajectory does not depend on the context. Rule $110$R has no such cell, of any length or kind. The proof is by hand.

\emph{Reach of the edge-move method} (Section~\ref{cl:sec:reach}). The proof of \eqref{imp:110} works at a stuck edge, a cell beyond which a configuration agrees with a spatially periodic background and which does not move under $F^{\pm1}$. In at least $12$ classes, among them $10$ open ones and the classes of $94$ and $95$, no configuration has a stuck edge; the proof is by hand (Proposition~\ref{cl:prop:nostuck}). For $11$ further open classes a background of small period with a stuck edge has no other fixed point of $\sh^aF^b$ next to it, for small $(a,b)$ (Proposition~\ref{cl:prop:reach}, computer-assisted), and for the remaining $7$ sampled computations indicate where the method stops (Remark~\ref{cl:rem:reach}).

\subsection{Context}
The problem of finding the block maps that commute with a given one goes back to Coven, Hedlund and Rhodes \cite{CHR1979}. Voorhees \cite{Voorhees1993} describes the binary cellular automata commuting with a given one as the solutions of a system of equations in the rule table, Moore and Boykett \cite{MooreBoykett1997} show that if one of two commuting cellular automata is permutive, the other can be expressed in terms of it, and Salo and T\"orm\"a \cite{SaloTorma2013} show that the commutant of a bipermutive automaton is small and that of a bipermutive affine one consists of affine automata. Boykett \cite{Boykett2004} lists by machine, up to isomorphism, the reversible one-dimensional cellular automata with a two-cell neighbourhood whose inverse also has a two-cell neighbourhood, for up to $12$ states. The censuses of this paper enumerate, for one reversible automaton and one radius at a time, all block maps on the four-letter alphabet that commute with it.

For rule $54$, by \eqref{imp:54}, the automorphism group of the second-order automaton is not virtually $\Z^d$, while for rule $110$ it is $\Z^2$ by \eqref{imp:110}. The first-order rules $110$ and $54$ are the two elementary rules that Wolfram singles out for class~4 behaviour, $110$ as the clear example and $54$ as a possible one \cite[pp.~693--697]{Wolfram2002}; their second-order versions thus have automorphism groups of opposite kinds, and the behaviour of the first-order rule is no guide to the automorphism group of $F_f$.
In $39$ of the $43$ classes of status (G) the room comes from a conserved local structure, a wall, with an autonomous interior; walls in each of these $39$ classes were exhibited by Takesue \cite[Section~4]{Takesue1989}, and what \cite{PaperC} adds is their consequence for the automorphism group. Of the other four, three are affine, and in the class of $95$ the room comes from the automorphisms of a system of light-speed particles \cite{PaperE}. The spacetime $X_F$ has zero $\Z^2$-entropy, so Hochman's theorem that positive entropy forces every finite group into the automorphism group of a $\Z^2$ shift of finite type \cite{Hochman2010} does not apply. The techniques that prove automorphism groups small do not apply either: the complexity arguments of Cyr--Kra and Donoso--Durand--Maass--Petite \cite{CK2015,DDMP2016} are for one-dimensional subshifts whose word complexity grows at most linearly or is not superlinear, while the rows of $X_F$ are all of $\A^n$; the algebraic rigidity of Kitchens and Schmidt \cite{KS1989,KS2000}, which gives $\Z^2$ for Ledrappier's shift \cite{Ledrappier1978} (\cite{Ward1994}, \cite[Theorem~7]{BRY2018}) and determines the groups of the affine classes of $15$, $60$, $105$ and $153$ \eqref{imp:affine}, needs a group shift. Ryan-type results \cite{Ryan1972,Salo2019Ryan,Kopra2022} compute centralizers of the whole group or of specially constructed sets of two or more automata, not of one given automaton. The arguments of the companion papers are specific to the rules; whether the property $\Aut(X_F)=\langle\sh,F\rangle$ of a reversible automaton is decidable is not known (Section~\ref{sec:discussion}).

\medskip
\noindent\textbf{Organisation.} Section~\ref{sec:prelim} fixes notation and defines blocking words, window graphs and eigenfields. Section~\ref{sec:method} describes the SAT encoding, the invertibility test, the verification and the cross-check, and Section~\ref{sec:110} gives the census for $110$ and $30$. Section~\ref{sec:classification} defines the statuses, recalls Theorem~\ref{cl:thm:A}, gives the tables and the proof of Theorem~\ref{thm:main}, states the census of the open classes, proves Proposition~\ref{cl:prop:nowalls}, records what is known about the open classes, and determines the reach of the edge-move method on them. Section~\ref{sec:discussion} states open questions. The appendix describes the computations.
